% !TeX root = Report.tex
% Introduction body. The chapter heading is declared in Report.tex.

Задача обнаружения прямых линий на изображении и оценивания их параметров возникает во многих задачах компьютерного зрения, анализа изображений и автоматической интерпретации данных. Даже в простой постановке, когда сигнал состоит из одной прямой на темном фоне, наличие шума существенно осложняет восстановление ее положения: наблюдаемое изображение содержит многочисленные шумовые выбросы, из-за которых положение прямой определяется менее устойчиво. В таких условиях требуется не только выделить линию визуально, но и надежно оценить ее параметры.

Одним из наиболее распространенных методов оценивания параметров прямой является преобразование Хафа \cite{duda1972hough}. Несмотря на свою эффективность, этот метод существенно зависит от качества входного изображения: аддитивный шум порождает пиксели, не принадлежащие сигналу, но по интенсивности схожие с ним, и тем самым ухудшает точность оценивания. Поэтому на практике преобразование Хафа обычно применяется не к исходному зашумленному изображению, а после предварительного шумоподавления.

В качестве такой предобработки естественно рассматривать классические методы шумоподавления, широко используемые в задачах обработки изображений, в частности медианный фильтр, фильтр Винера и метод \texttt{BM3D} \cite{huang1979,wiener1949,lim1990,dabov2007}. Наряду с ними представляет интерес более специальный алгоритм предобработки \cite{trickett2003}, предложенный С. Трикеттом в 2003 году. Теоретической основой этого алгоритма служит метод анализа сингулярного спектра SSA и его комплекснозначный вариант Complex SSA (CSSA) \cite{ts_structure,ssa_with_r,zhornikova2016}. Для изображений это приводит к схеме, в которой к строкам или столбцам матрицы после преобразования Фурье применяется CSSA, а затем выполняется обратное преобразование Фурье. Таким образом, в данной работе CSSA рассматривается как часть другого алгоритма.

В данной работе исследуется именно такая постановка: различные методы шумоподавления сравниваются по тому, как они влияют на точность последующего оценивания параметров прямой методом Хафа.

Рассматривается модель наблюдения
\[
\mathbf{M} = \mathbf{S} + \mathbf{R},
\]
где $\mathbf{M} \in \mathbb{R}^{m \times n}$ --- наблюдаемое изображение, $\mathbf{S}$ --- компонента сигнала, а шумовая матрица $\mathbf{R}$ имеет независимые элементы, распределенные по нормальному закону $\mathcal{N}(0,\sigma^2)$. В данной работе предполагается, что компонента сигнала устроена следующим образом: пикселям, лежащим на прямой, соответствует фиксированная положительная интенсивность, а вне прямой значения равны нулю. В численных экспериментах основное внимание уделяется случаю одной прямой, задаваемой уравнением $y = ax + b$. Итоговое сравнение проводится по точности оценивания параметров $\rho$ и $\theta$.

В качестве базовых методов предобработки рассматриваются три классических метода: медианный фильтр, фильтр Винера и \texttt{BM3D}. С ними сравнивается алгоритм предобработки на основе \texttt{CSSA row.row}. Далее анализируется, насколько такой алгоритм предобработки может конкурировать с классическими фильтрами по точности последующего оценивания параметров прямой методом Хафа.
% TODO: завершить формулировку центрального вопроса работы.

\textbf{Цель работы} состоит в сравнении различных методов предобработки изображения по точности оценивания параметров прямой методом Хафа.

Для достижения этой цели в работе решаются следующие задачи:
\begin{enumerate}
	\item описать используемые алгоритмы: преобразование Хафа, классические методы шумоподавления и алгоритм предобработки, основанный на Complex SSA;
	\item исследовать работу варианта \texttt{CSSA row.row} в задаче подавления шума для одной прямой, включая выбор числа компонент, влияние расположения прямой и уровня шума, а также типичные ошибки восстановления;
	\item провести численное сравнение методов предобработки по точности оценок параметров $\rho$ и $\theta$ после применения преобразования Хафа;
	\item интерпретировать полученные результаты, опираясь на свойства и ограничения метода CSSA, и наметить дальнейшие направления исследования.
\end{enumerate}

Кратко опишем структуру работы. В первой главе описываются используемые алгоритмы и вводятся обозначения, необходимые для дальнейшего изложения. Во второй главе подробно исследуется использование \texttt{CSSA row.row} для подавления шума в случае одной прямой и обсуждаются ограничения этого подхода. В третьей главе приводится численное сравнение различных предобработок по точности оценивания параметров $\rho$ и $\theta$. В заключении подводятся итоги исследования и обсуждаются возможные направления дальнейшей работы, в том числе автоматический выбор числа компонент, автоматическое определение информативной компоненты сигнала и переход к случаю нескольких прямых на изображении.
